% Resumo em Língua Inglesa
\setlength{\absparsep}{18pt}
\begin{resumo}[ABSTRACT]
 \begin{otherlanguage*}{english}

This work presents the development and evaluation of a gestural musical hyperinstrument named Aerochord, which utilizes computer vision with a standard RGB camera to generate MIDI 2.0 events with extended resolution, aiming to democratize musical expression for users without technical training on instruments and for people with reduced mobility. Initially, a literature review was conducted on human-computer interaction and computer music, followed by a critical analysis of the state of the art in gestural interfaces, which highlighted a gap between accessibility, latency, and professional-level expressivity. Based on this investigation, the functional and non-functional requirements for Aerochord were defined, including real-time video processing, hand landmark detection via computer vision techniques, gesture-to-sound mappings, and communication via MIDI 2.0 with a contingency plan for legacy mode. A modular architecture in C++ using the JUCE framework was implemented, organized into pipelines that minimize jitter, utilizing lock-free queues, handshaking via MIDI Capability Inquiry, and the packaging of commands into Universal MIDI Packets to ensure a coherent sonic response to the performer's gestures. The experimental evaluation carried out demonstrated end-to-end latency consistently below the perceptible threshold of 30 ms (95th percentile of 24.45 ms) and low resource consumption, alongside a pilot usability test that evidenced a low entry barrier for users without musical training. The results confirm the technical and practical feasibility of the proposal, establishing a foundation for future developments in polyphony, adaptive mappings, and immersive environments. It is expected that Aerochord will contribute to the fields of human-computer interaction and computer music by presenting a model of an accessible, responsive, and expressive hyperinstrument.


% Coloque após os ":"  as palavras-chave com a primeira 
% letra de cada palavra em maiúscula e separadas por ponto.
\textbf{Keywords}: gesture interfaces, musical computing, MIDI 2.0, musical democratization, computer vision, hyperinstruments, accessibility, low latency.
 \end{otherlanguage*}
\end{resumo}